\chapter{Sprint 0 : Preparation et analyse}

\section{Introduction}
Tout projet Scrum debute par ce que l'on appelle le Sprint 0 : une phase
de preparation dont l'objectif n'est pas de livrer du code mais d'eliminer
les incertitudes sur \emph{ce} qui va etre livre. Le Sprint 0 est, dans
notre cas, l'etape ou nous avons identifie les acteurs interagissant avec
la plateforme, redige les besoins fonctionnels et non fonctionnels qui
allaient guider la priorisation, fixe l'environnement technique, et
esquisse les diagrammes d'analyse et de conception sur lesquels la mise
en oeuvre se referera ulterieurement.

Ce chapitre est organise en cinq sections. Nous commencons par les acteurs.
Nous decoupons ensuite les besoins en categories fonctionnelle et non
fonctionnelle. Nous decrivons la pile technique --- back-end, front-end,
fournisseurs d'IA, hebergement --- en explicitant le raisonnement derriere
chaque choix majeur. Nous presentons ensuite les artefacts de conception
(diagramme de cas d'utilisation, modele conceptuel des donnees, diagramme
de sequence pour le parcours de prise de rendez-vous), et nous cloturons
par une vue d'ensemble de la planification reliant chaque sprint suivant
aux exigences qu'il satisfera.

\section{Identification des acteurs}
Quatre acteurs principaux interagissent avec Nova~AI. Chacun dispose d'un
point d'entree distinct dans le systeme et d'attentes specifiques.

\subsection{Le proprietaire d'entreprise}
Le proprietaire est l'utilisateur payant. Dans une clinique mono-medecin,
il est aussi le medecin. Dans une configuration multi-praticiens, il est
typiquement un medecin senior ou un gestionnaire administratif. Dans un
laboratoire, il est le biologiste responsable. Dans un garage, il est le
chef mecanicien.

Le proprietaire entre dans le systeme en s'inscrivant sur la page
d'accueil marketing, en choisissant un type d'entreprise dans la liste
deroulante (clinique / laboratoire / cabinet medical / salon / garage /
entreprise), puis en completant le flux de creation d'entreprise. Une
fois entre, il configure les services, les membres de l'equipe, les
heures de travail, les forfaits et les diverses fonctionnalites IA. Il
consulte les tableaux de bord, approuve les campagnes de remplissage de
creneaux, valide les surcharges de securite des prescriptions, examine
les suggestions d'apprentissage de duree de consultation, scanne les
codes QR des patients pour demarrer les consultations.

\subsection{Le praticien (membre de l'equipe)}
Le praticien est un membre de l'equipe : un medecin dans une clinique,
un biologiste dans un laboratoire, un mecanicien dans un garage, un
coiffeur dans un salon. Il est distinct du proprietaire car ses
permissions sont plus restreintes : il devrait voir sa journee, ses
patients, son historique de rendez-vous, mais pas les tableaux de bord
financiers, les surfaces de gestion d'equipe ou les analyses globales.

Le praticien entre soit par la sous-route \emph{Vue medecin} du tableau
de bord (lorsque le proprietaire clique sur \og Voir comme \fg{} sur un
membre d'equipe) soit par la PWA patient (lorsque le membre d'equipe a
ete lie a un \texttt{auth\_uid} et se connecte directement). Cette
deuxieme voie est plus courante dans les configurations multi-praticiens
ou chaque medecin dispose de ses propres identifiants.

\subsection{Le patient ou client}
Le patient (ou client, dans les verticaux non medicaux) est l'utilisateur
final. Il interagit avec Nova principalement via WhatsApp : il envoie des
demandes de rendez-vous, pose des questions FAQ, envoie des photos de
prescriptions ou de pieces de voiture endommagees, repond aux rappels.
Certains patients vont plus loin et installent la PWA sur leur telephone,
ce qui leur donne acces a une experience plus riche : position dans la
file d'attente avec ETA en temps reel, check-in QR a la clinique,
chronologie de fichiers, profil.

Crucialement, le patient n'a jamais besoin de s'inscrire pour utiliser
le cote WhatsApp de la plateforme. La premiere fois qu'il envoie un
message a une entreprise pilotee par Nova, une ligne est automatiquement
creee dans la table \texttt{clients} et il devient une entite suivie a
partir de ce moment.

\subsection{Le systeme (agents automatises)}
Le quatrieme acteur est le systeme lui-meme. Un nombre croissant de
fonctionnalites s'executent selon une planification plutot qu'en reponse
a une action utilisateur : les rappels 24~heures et 1~heure declenches
par \texttt{send-reminders}, le scan quotidien de remplissage de
creneaux, le lot quotidien d'anniversaire qui frappe les vouchers, la
revue hebdomadaire d'apprentissage de duree qui ajuste les durees de
service en fonction de la realite observee, et la campagne hebdomadaire
de reactivation pour les clients dormants.

Ces agents automatises sont des citoyens de premiere classe dans le
diagramme de cas d'utilisation, car leur comportement fait partie de la
proposition de valeur : un proprietaire de clinique achete Nova en
partie parce que Nova fait des choses \emph{pendant que le proprietaire
dort}.

\section{Specification des besoins}

\subsection{Besoins fonctionnels}
Nous avons catalogue les besoins fonctionnels dans un tableau unique
organise par acteur (Tableau \ref{tab:fr}). Le tableau a evolue
progressivement au fil des sprints --- le Sprint 1 a couvert les vingt
premieres lignes, le Sprint 2 en a ajoute quinze, le Sprint 3 a cloture
la liste avec les elements pilotes par IA.

\begin{table}[H]
  \centering\footnotesize
  \begin{tabular}{p{2.4cm}p{1.2cm}p{9.8cm}}
    \toprule
    \textbf{Acteur} & \textbf{ID} & \textbf{Besoin} \\
    \midrule
    Proprietaire & BF01 & S'inscrire avec email et lien magique, puis creer son profil d'entreprise. \\
    Proprietaire & BF02 & Choisir un type d'entreprise dans la liste supportee, avec valeurs par defaut sectorielles pre-chargees. \\
    Proprietaire & BF03 & Ajouter et gerer des services (nom, duree, prix, categorie). \\
    Proprietaire & BF04 & Ajouter et gerer les membres de l'equipe (medecins, mecaniciens, coiffeurs). \\
    Proprietaire & BF05 & Configurer le receptionniste IA : activation, heures d'ouverture, FAQ, jours travailles. \\
    Proprietaire & BF06 & Visualiser le calendrier journalier, hebdomadaire et mensuel, filtrable par membre d'equipe. \\
    Proprietaire & BF07 & Creer, modifier et annuler manuellement des rendez-vous. \\
    Proprietaire & BF08 & Consulter l'historique complet de chaque client (rendez-vous, messages, fichiers, resultats). \\
    Proprietaire & BF09 & Recevoir des notifications navigateur en temps reel pour nouvelles reservations, triages urgents, no-shows. \\
    Proprietaire & BF10 & Configurer les campagnes de vouchers d'anniversaire et consulter l'historique. \\
    Proprietaire & BF11 & Approuver les campagnes de remplissage de creneaux proposees par le scanner quotidien. \\
    Proprietaire & BF12 & Activer l'auto-ajustement de duree par service et examiner les suggestions de derive. \\
    Proprietaire & BF13 & Scanner les codes QR patients depuis la barre du tableau de bord pour demarrer les consultations. \\
    Proprietaire & BF14 & Lancer des verifications IA de securite des prescriptions contre l'historique du patient. \\
    Proprietaire & BF15 & Gerer les vehicules (entreprises de type garage) : VIN, plaque, kilometrage, photos, journal de service. \\
    Proprietaire & BF16 & Generer et imprimer le code QR de reception (laboratoires) pour la file d'attente en temps reel. \\
    Praticien & BF17 & Se connecter a la PWA et voir les rendez-vous de la journee. \\
    Praticien & BF18 & Marquer un patient comme Arrive et Termine (horodatage debut/fin consultation). \\
    Praticien & BF19 & Scanner un QR patient avec la camera du telephone pour ouvrir le minuteur de consultation. \\
    Praticien & BF20 & Cliquer sur Terminer Consultation a la fin de la visite, enregistrant la duree reelle. \\
    Patient & BF21 & Envoyer un message WhatsApp dans l'une des trois langues et recevoir une reponse pertinente. \\
    Patient & BF22 & Reserver un rendez-vous via une conversation multi-tour (nom, age, service, date, heure). \\
    Patient & BF23 & Confirmer ou annuler un rendez-vous rappele via les boutons interactifs WhatsApp. \\
    Patient & BF24 & Envoyer une photo d'une prescription manuscrite et recevoir une liste extraite et corrigee des analyses. \\
    Patient & BF25 & Envoyer une photo de degats voiture et la voir auto-attachee a son fichier vehicule. \\
    Patient & BF26 & Scanner le QR de reception (laboratoire) et rejoindre la file d'attente en temps reel. \\
    Patient & BF27 & Voir sa position dans la file et le temps d'attente estime, mis a jour en direct. \\
    Patient & BF28 & Recevoir un WhatsApp personnalise le jour de son anniversaire avec un code voucher. \\
    Patient & BF29 & Installer la PWA Nova sur son telephone et la consulter hors ligne. \\
    Patient & BF30 & Verifier son numero de telephone via un OTP a six chiffres. \\
    Systeme & BF31 & Envoyer les rappels 24h pour rendez-vous confirmes, avec boutons Confirmer/Annuler. \\
    Systeme & BF32 & Envoyer les rappels 1h pour rendez-vous confirmes. \\
    Systeme & BF33 & Detecter chaque matin les opportunites de creneaux vides et creer des brouillons de campagne. \\
    Systeme & BF34 & Executer le batch d'anniversaire chaque matin, frapper les codes, envoyer les WhatsApp. \\
    Systeme & BF35 & Recalculer les moyennes de duree de service hebdomadairement et produire suggestions de derive. \\
    Systeme & BF36 & Reactiver les clients dormants (derniere visite > 6 mois) avec une campagne hebdomadaire. \\
    Systeme & BF37 & Auto-redemption d'un voucher d'anniversaire lorsque le client envoie le code. \\
    Systeme & BF38 & Auto-attacher les photos garage au vehicule correspondant en analysant le message. \\
    \bottomrule
  \end{tabular}
  \caption{Catalogue des besoins fonctionnels}
  \label{tab:fr}
\end{table}

\subsection{Besoins non fonctionnels}
Les besoins non fonctionnels sont tout aussi importants. Certains, comme
l'isolation multi-locataire, n'etaient pas negociables car leur defaillance
serait catastrophique. D'autres, comme la cible des cinq secondes de
reponse, etaient plus souples mais ont neanmoins guide les choix
architecturaux.

\begin{itemize}
  \item \textbf{Multi-tenancy avec isolation stricte.} Chaque entreprise
        ne voit que ses propres donnees, applique au niveau de la base
        via les politiques Row-Level Security de Postgres.
  \item \textbf{Mises a jour en temps reel.} Les tuiles du tableau de
        bord, les positions de file d'attente et les minuteurs de
        consultation se mettent a jour en moins de deux secondes via
        Supabase Realtime.
  \item \textbf{Support multilingue.} Le systeme gere l'arabe dialectal,
        le francais et l'anglais en messages entrants, reponses
        sortantes, gabarits de vouchers et textes de rappel.
  \item \textbf{Mobile-first.} La PWA est entierement fonctionnelle sur
        telephones jusqu'a 360 pixels de largeur ; le tableau de bord
        redirige automatiquement les proprietaires mobiles connectes
        vers la PWA.
  \item \textbf{Fenetre de reponse de cinq secondes.} L'IA doit repondre
        en moins de cinq secondes pour 95\% des cas. Pour atteindre cet
        objectif, la fonction Edge utilise Claude Haiku (plus rapide
        que Sonnet) et applique un timeout strict de six secondes via
        \texttt{Promise.race}.
  \item \textbf{Tolerance aux pannes.} Un appel IA defaillant ne doit
        jamais bloquer le flux de prise de rendez-vous. Le message
        entrant est insere d'abord, puis l'appel IA est tente.
  \item \textbf{Conformite RGPD.} Les numeros de telephone et
        conversations sont attaches a l'entreprise concernee et
        supprimables sur demande.
  \item \textbf{Maitrise des couts.} Le cout total d'infrastructure par
        locataire et par mois doit rester sous 5~USD au regime de
        croisiere, pour qu'un abonnement mensuel a 50~TND soit rentable.
\end{itemize}

\section{Environnement technique}
\subsection{Pile back-end}
Le back-end est presque entierement construit sur Supabase, avec quelques
fournisseurs d'IA externes attaches. Le tableau \ref{tab:back} resume les
choix.

\begin{table}[H]
  \centering\footnotesize
  \begin{tabular}{p{3.5cm}p{4cm}p{6cm}}
    \toprule
    \textbf{Composant} & \textbf{Choix} & \textbf{Justification} \\
    \midrule
    Base de donnees & Postgres 15 (Supabase) & Mature, JSONB pour colonnes flexibles, RLS pour multi-tenancy. \\
    Auth & Supabase Auth (lien magique, OAuth) & Integre nativement avec RLS via \texttt{auth.uid()}. \\
    Fonctions & Supabase Edge Functions (Deno) & Distribuees globalement, V8 isolates, faible cold-start. \\
    Realtime & Supabase Realtime & Pousse les changements Postgres aux clients via WebSockets. \\
    Storage & Supabase Storage & Compatible S3, RLS, URLs signees. \\
    Cron & pg\_cron + pg\_net & Planification dans la base sans worker externe. \\
    Classifieur d'intention & Anthropic Claude Haiku 4.5 & Meilleur ratio vitesse/cout pour JSON structure court. \\
    Vision OCR & Anthropic Claude Sonnet 4.6 & Robuste sur ecriture francaise et abreviations medicales. \\
    Speech-to-text & OpenAI Whisper & Robuste aux notes vocales bruyantes. \\
    Canal conversationnel & WhatsApp Business Cloud API & Direct depuis Meta, sans agregateur. \\
    \bottomrule
  \end{tabular}
  \caption{Pile back-end}
  \label{tab:back}
\end{table}

\subsection{Pile front-end}
Le front-end fait des choix tout aussi opinionnes : JavaScript ES Modules
vanille (sans framework) pour le tableau de bord et la PWA, CSS fait main
avec un theme \emph{anti-gravite} personnalise (orbes flottantes, cartes
qui levitent, gradient violet-cyan), et hebergement statique sur Vercel
avec deploiement automatique a chaque push GitHub.

\subsection{Pourquoi pas de framework}
Le tableau de bord est riche en fonctionnalites mais sa complexite est
bornee. Un framework moderne aurait impose une etape de build, un bundler,
un pipeline de deploiement, et un cout de base non trivial de JavaScript
livre a l'utilisateur. En restant sur des modules ES vanille servis depuis
un seul fichier HTML, nous gardons l'ensemble du projet sous 350~Ko
gzippes, deployons en moins de trente secondes, et evitons toute une
classe de problemes lies aux dependances.

\section{Diagramme de cas d'utilisation}
% ===========================================================================
% Diagramme de cas d'utilisation global — Nova AI
% ===========================================================================
\begin{figure}[H]
\centering
\begin{tikzpicture}[
  font=\footnotesize,
  every node/.style={align=center},
  actor/.style ={font=\footnotesize\bfseries},
  uc/.style    ={ellipse, draw=black!75, fill=cyan!8,
                 minimum width=32mm, minimum height=8mm,
                 font=\scriptsize, inner sep=2pt},
  inc/.style   ={->, thick, draw=black!70, dashed, >=Stealth},
  rel/.style   ={-, thick, draw=black!75},
  sys/.style   ={rectangle, draw=black!60, fill=white,
                 minimum width=88mm, minimum height=110mm,
                 line width=1pt}
]

% System box (Nova AI)
\node[sys] (sys) at (0,0) {};
\node[font=\bfseries, anchor=north] at (sys.north) {\bfseries Système Nova~AI};

% Use-case nodes inside the system, arranged in two columns
% Left column — patient flow
\node[uc] (book)     at (-2.0, 4.2) {Réserver un\\rendez-vous};
\node[uc] (cancel)   at (-2.0, 3.0) {Annuler un\\rendez-vous};
\node[uc] (resched)  at (-2.0, 1.8) {Reprogrammer\\un rendez-vous};
\node[uc] (uploadrx) at (-2.0, 0.6) {Envoyer une\\ordonnance};
\node[uc] (askinfo)  at (-2.0,-0.6) {Demander des\\informations};
\node[uc] (urgent)   at (-2.0,-1.8) {Signaler une\\urgence};

% Right column — staff flow
\node[uc] (viewcal)  at ( 2.0, 4.2) {Consulter le\\calendrier};
\node[uc] (mngsvc)   at ( 2.0, 3.0) {Gérer les\\services};
\node[uc] (mngteam)  at ( 2.0, 1.8) {Gérer l'équipe};
\node[uc] (revenue)  at ( 2.0, 0.6) {Voir les\\indicateurs};
\node[uc] (asknova)  at ( 2.0,-0.6) {Interroger Nova\\(vocal)};
\node[uc] (settings) at ( 2.0,-1.8) {Configurer\\horaires \& IA};

% Shared
\node[uc] (auto-ai)  at ( 0.0,-3.4) {Classifier l'intention\\(IA)};

% Actors outside the system
\node[actor]                                 (patient) at (-7.5, 1.5) {\faUser\ Patient};
\node[actor, font=\footnotesize\bfseries]    (owner)   at ( 7.5, 2.5) {Propriétaire\\(médecin / gérant)};
\node[actor, font=\footnotesize\bfseries]    (recept)  at ( 7.5,-0.5) {Réceptionniste\\/ assistant};
\node[actor, font=\footnotesize\bfseries]    (system)  at ( 0.0,-5.0) {\textit{<<système IA>>}\\Anthropic Claude};

% Patient relationships
\draw[rel] (patient) -- (book);
\draw[rel] (patient) -- (cancel);
\draw[rel] (patient) -- (resched);
\draw[rel] (patient) -- (uploadrx);
\draw[rel] (patient) -- (askinfo);
\draw[rel] (patient) -- (urgent);

% Owner relationships
\draw[rel] (owner) -- (viewcal);
\draw[rel] (owner) -- (mngsvc);
\draw[rel] (owner) -- (mngteam);
\draw[rel] (owner) -- (revenue);
\draw[rel] (owner) -- (asknova);
\draw[rel] (owner) -- (settings);

% Receptionist subset
\draw[rel] (recept) -- (book);
\draw[rel] (recept) -- (resched);
\draw[rel] (recept) -- (viewcal);

% AI inclusion
\draw[inc] (book)    -- node[midway, sloped, above=-1pt, font=\tiny]{<<include>>} (auto-ai);
\draw[inc] (askinfo) -- node[midway, sloped, above=-1pt, font=\tiny]{<<include>>} (auto-ai);
\draw[inc] (urgent)  -- node[midway, sloped, above=-1pt, font=\tiny]{<<include>>} (auto-ai);
\draw[inc] (asknova) -- node[midway, sloped, above=-1pt, font=\tiny]{<<include>>} (auto-ai);
\draw[rel] (system)  -- (auto-ai);

\end{tikzpicture}
\caption{Diagramme de cas d'utilisation global : trois acteurs humains (patient, propriétaire, réceptionniste) interagissent avec le système, qui s'appuie sur l'acteur secondaire \textit{<<système IA>>} Anthropic Claude pour classifier les intentions et conduire les conversations multi-tours.}
\label{fig:uc}
\end{figure}


\section{Modele conceptuel des donnees}
\begin{figure}[!htbp]\centering
\includegraphics[width=.98\textwidth,height=.85\textheight,keepaspectratio]{figures/erd.png}
\caption{Modèle conceptuel des données.}\label{fig:erd}
\end{figure}


\section{Diagramme de sequence du parcours de prise de rendez-vous}
% ===========================================================================
% Diagramme de séquence — parcours de prise de rendez-vous multi-tour
% ===========================================================================
\begin{figure}[H]
\centering
\begin{tikzpicture}[
  font=\scriptsize,
  every node/.style={align=center},
  actor/.style ={rectangle, rounded corners=2pt, draw=black!75,
                 fill=cyan!12, minimum width=22mm, minimum height=8mm,
                 font=\scriptsize\bfseries},
  msg/.style   ={->, thick, draw=black!80, >=Stealth},
  msg-d/.style ={->, thick, draw=black!60, >=Stealth, dashed},
  ret/.style   ={<-, thick, draw=black!60, >=Stealth, dashed},
  lab/.style   ={font=\scriptsize, fill=white, inner sep=1pt},
]

% Vertical lifelines
\foreach \i/\name/\x in {0/Patient/0, 1/{Meta WhatsApp}/3, 2/{whatsapp-\\webhook}/6, 3/{classify-\\message}/9, 4/{Anthropic\\Claude}/12, 5/{Postgres\\(Supabase)}/15} {
  \node[actor] (a\i) at (\x, 0) {\name};
  \draw[dashed, gray!60] (a\i.south) -- ++(0,-13);
}

% Step 1 — patient sends "نحب نعمل ECG"
\node[lab] at (1.5, -1)   {1. Message WhatsApp \og je veux un ECG \fg{}};
\draw[msg] (a0.south)|-(1.5, -1)-|(a1.north);

% Step 2 — Meta delivers webhook
\node[lab] at (4.5, -2.0) {2. POST webhook (texte + métadonnées)};
\draw[msg] (a1.south)|-(4.5, -2.0)-|(a2.north);

% Step 3 — webhook stores message + triggers classify
\node[lab] at (7.5, -3.0) {3. INSERT messages, appel classify};
\draw[msg] (a2.south)|-(7.5, -3.0)-|(a3.north);

% Step 4 — classify calls Claude for intent + parsing
\node[lab] at (10.5, -4.0) {4. Prompt strict-JSON};
\draw[msg] (a3.south)|-(10.5, -4.0)-|(a4.north);
\node[lab] at (10.5, -4.7) {5. \texttt{intent="book"} + champs};
\draw[ret] (a3.south)|-(10.5, -4.7)-|(a4.north);

% Step 5 — classify writes booking_drafts
\node[lab] at (12, -5.6)   {6. UPSERT booking\_drafts (state="time")};
\draw[msg] (a3.south)|-(12, -5.6)-|(a5.north);

% Step 6 — classify formats clarification + sends back
\node[lab] at (4.5, -6.6)  {7. Réponse \og Quelle heure ? \fg{} via API Meta};
\draw[ret] (a2.south)|-(4.5, -6.6)-|(a3.north);
\draw[msg] (a2.south)|-(1.5, -7.4)-|(a1.north);
\node[lab] at (1.5, -7.4)  {8. Question reçue par patient};
\draw[msg] (a1.south)|-(1.5, -8.2)-|(a0.north);

% Step 7 — patient replies "14h"
\node[lab] at (1.5, -9.2)  {9. Réponse \og 14h \fg{}};
\draw[msg] (a0.south)|-(1.5, -9.2)-|(a1.north);
\draw[msg] (a1.south)|-(4.5, -9.9)-|(a2.north);
\draw[msg] (a2.south)|-(7.5, -10.4)-|(a3.north);

% Step 8 — classify validates working_hours + buffer + creates appointment
\node[lab] at (12, -11.2)  {10. SELECT working\_hours \& slots};
\draw[msg] (a3.south)|-(12, -11.2)-|(a5.north);
\node[lab] at (12, -11.9)  {11. INSERT appointments status="confirmed"};
\draw[msg] (a3.south)|-(12, -11.9)-|(a5.north);

% Step 9 — confirmation back
\node[lab] at (4.5, -12.6) {12. \og C'est confirmé !\fg{} via API Meta};
\draw[ret] (a2.south)|-(4.5, -12.6)-|(a3.north);

\end{tikzpicture}
\caption{Diagramme de séquence simplifié de la prise de rendez-vous multi-tours : du message WhatsApp initial du patient jusqu'à la création effective de la ligne \texttt{appointments} et le retour de confirmation. Les flèches pointillées indiquent les réponses asynchrones.}
\label{fig:seq}
\end{figure}


\section{Planification des sprints}
Les quatre sprints ont ete planifies comme suit. Chaque sprint a un objectif
clair, une duree fixe de deux a trois semaines, et un backlog tire du
tableau des besoins.

\begin{table}[H]
  \centering\footnotesize
  \begin{tabular}{p{1.6cm}p{2.6cm}p{9cm}}
    \toprule
    \textbf{Sprint} & \textbf{Objectif} & \textbf{Faits saillants du backlog} \\
    \midrule
    Sprint 0 & Preparation & Acteurs, besoins, pile technique, diagrammes de conception. \\
    Sprint 1 & Fondations & Schema, auth, RLS, webhook WhatsApp, classifieur IA, prise de rendez-vous multi-tour. \\
    Sprint 2 & Operations & Tableau de bord, calendrier, rendez-vous, pipeline labo, OCR de prescriptions, PWA, OTP telephone. \\
    Sprint 3 & Intelligence & Demand-fill, apprentissage des durees, securite des medicaments, profils vehicules, anniversaires, check-in QR, file d'attente. \\
    \bottomrule
  \end{tabular}
  \caption{Vue d'ensemble de la planification des sprints}
  \label{tab:plan}
\end{table}

\section{Conclusion}
Le Sprint 0 a pose le decor. Nous avons identifie les quatre acteurs,
catalogue trente-huit besoins fonctionnels et neuf besoins non
fonctionnels, justifie les choix techniques majeurs, et esquisse les
artefacts d'analyse et de conception qui guideront la mise en oeuvre.
Avec ces fondations en place, le chapitre suivant entre dans le Sprint 1,
ou la colonne vertebrale multi-locataire de la plateforme a ete construite.
